\documentclass[10pt,a4paper]{amsart}
\usepackage{amsmath,amssymb,amsthm,graphicx}
% For commutative diagrams we use Xy-pic
\usepackage[all]{xy}

% We use 2cell for 2-commutative diagrams.
\xyoption{2cell}
% Theorem environments.
%
\theoremstyle{plain}
\newtheorem{theorem}[subsection]{Theorem}
\newtheorem{proposition}[subsection]{Proposition}
\newtheorem{lemma}[subsection]{Lemma}

\theoremstyle{definition}
\newtheorem{definition}[subsection]{Definition}
\newtheorem{example}[subsection]{Example}
\newtheorem{exercise}[subsection]{Exercise}
\newtheorem{situation}[subsection]{Situation}

\theoremstyle{remark}
\newtheorem{remark}[subsection]{Remark}
\newtheorem{remarks}[subsection]{Remarks}

\numberwithin{equation}{subsection}

% Macros
%
\def\lim{\mathop{\mathrm{lim}}\nolimits}
\def\colim{\mathop{\mathrm{colim}}\nolimits}
\def\Spec{\mathop{\mathrm{Spec}}}
\def\Hom{\mathop{\mathrm{Hom}}\nolimits}
\def\Ext{\mathop{\mathrm{Ext}}\nolimits}
\def\SheafHom{\mathop{\mathcal{H}\!\mathit{om}}\nolimits}
\def\SheafExt{\mathop{\mathcal{E}\!\mathit{xt}}\nolimits}
\def\Sch{\mathit{Sch}}
\def\Mor{\mathop{\mathrm{Mor}}\nolimits}
\def\Ob{\mathop{\mathrm{Ob}}\nolimits}
\def\Sh{\mathop{\mathit{Sh}}\nolimits}
\def\NL{\mathop{N\!L}\nolimits}
\def\CH{\mathop{\mathrm{CH}}\nolimits}
\def\proetale{{pro\text{-}\acute{e}tale}}
\def\etale{{\acute{e}tale}}
\def\QCoh{\mathit{QCoh}}
\def\Ker{\mathop{\mathrm{Ker}}}
\def\Im{\mathop{\mathrm{Im}}}
\def\Coker{\mathop{\mathrm{Coker}}}
\def\Coim{\mathop{\mathrm{Coim}}}

% Boxtimes
%
\DeclareMathSymbol{\boxtimes}{\mathbin}{AMSa}{"02}

%
% Macros for moduli stacks/spaces
%
\def\QCohstack{\mathcal{QC}\!\mathit{oh}}
\def\Cohstack{\mathcal{C}\!\mathit{oh}}
\def\Spacesstack{\mathcal{S}\!\mathit{paces}}
\def\Quotfunctor{\mathrm{Quot}}
\def\Hilbfunctor{\mathrm{Hilb}}
\def\Curvesstack{\mathcal{C}\!\mathit{urves}}
\def\Polarizedstack{\mathcal{P}\!\mathit{olarized}}
\def\Complexesstack{\mathcal{C}\!\mathit{omplexes}}
% \Pic is the operator that assigns to X its picard group, usage \Pic(X)
% \Picardstack_{X/B} denotes the Picard stack of X over B
% \Picardfunctor_{X/B} denotes the Picard functor of X over B
\def\Pic{\mathop{\mathrm{Pic}}\nolimits}
\def\Picardstack{\mathcal{P}\!\mathit{ic}}
\def\Picardfunctor{\mathrm{Pic}}
\def\Deformationcategory{\mathcal{D}\!\mathit{ef}}
\newenvironment{reference}{}{}
\newenvironment{slogan}{}{}
\newenvironment{history}{}{}
\usepackage{cleveref}
\crefname{originaldefinition}{original definition}{original definitions}
\Crefname{originaldefinition}{Original definition}{Original definitions}
\makeatletter
\newlabel{definition-type}{{0C6Z}{source}{}{sourcepointer-1}{}}
\newlabel{definition-type@cref}{{[originaldefinition][0][]0C6Z}{[1][0][]source}{}{}{}}
\makeatother
\crefname{originaldefinition}{original definition}{original definitions}
\Crefname{originaldefinition}{Original definition}{Original definitions}
\makeatletter
\newlabel{definition-genus}{{0C72}{source}{}{sourcepointer-2}{}}
\newlabel{definition-genus@cref}{{[originaldefinition][0][]0C72}{[1][0][]source}{}{}{}}
\makeatother
\crefname{originalequation}{original equation}{original equations}
\Crefname{originalequation}{Original equation}{Original equations}
\makeatletter
\newlabel{equation-form}{{0C65}{source}{}{sourcepointer-3}{}}
\newlabel{equation-form@cref}{{[originalequation][0][]0C65}{[1][0][]source}{}{}{}}
\makeatother
\crefname{originallemma}{original lemma}{original lemmas}
\Crefname{originallemma}{Original lemma}{Original lemmas}
\makeatletter
\newlabel{lemma-intersection-pairing}{{0C64}{source}{}{sourcepointer-4}{}}
\newlabel{lemma-intersection-pairing@cref}{{[originallemma][0][]0C64}{[1][0][]source}{}{}{}}
\makeatother
\crefname{originallemma}{original lemma}{original lemmas}
\Crefname{originallemma}{Original lemma}{Original lemmas}
\makeatletter
\newlabel{lemma-properties-form}{{0C66}{source}{}{sourcepointer-5}{}}
\newlabel{lemma-properties-form@cref}{{[originallemma][0][]0C66}{[1][0][]source}{}{}{}}
\makeatother
\crefname{originallemma}{original lemma}{original lemmas}
\Crefname{originallemma}{Original lemma}{Original lemmas}
\makeatletter
\newlabel{lemma-genus-formula}{{0CA3}{source}{}{sourcepointer-6}{}}
\newlabel{lemma-genus-formula@cref}{{[originallemma][0][]0CA3}{[1][0][]source}{}{}{}}
\makeatother
\crefname{originallemma}{original lemma}{original lemmas}
\Crefname{originallemma}{Original lemma}{Original lemmas}
\makeatletter
\newlabel{varieties-lemma-divisible}{{0C6P}{source}{}{sourcepointer-7}{}}
\newlabel{varieties-lemma-divisible@cref}{{[originallemma][0][]0C6P}{[1][0][]source}{}{}{}}
\makeatother
\crefname{originallemma}{original lemma}{original lemmas}
\Crefname{originallemma}{Original lemma}{Original lemmas}
\makeatletter
\newlabel{lemma-divisor-special-fiber}{{0C5Z}{source}{}{sourcepointer-8}{}}
\newlabel{lemma-divisor-special-fiber@cref}{{[originallemma][0][]0C5Z}{[1][0][]source}{}{}{}}
\makeatother
\title{The data read off a regular model form a numerical type whose genus is the genus of the curve}
\author{The Stacks Project Authors}
\date{stacks-project:models.tex:local-snapshot-20261012; GNU FDL 1.2 or later}
\begin{document}
\maketitle
\noindent\textit{Editorial note.} The counted unit is the Stacks Project Lemma lemma-numerical-type-of-model at source lines 5221-5230 together with its proof at source lines 5232-5264, a straight adjacent pairing in the chapter file. The chapter is the upstream file models.tex; the pool directory carries the longer name stacks-models and is not used as the chapter name anywhere in this card.\par\medskip
\noindent\textit{Editorial note.} One context span carries the regular-model situation at lines 4628-4637, which the retained statement names; the situation carries a reference of its own to the divisor lemma, which is therefore declared as a pointer as well. The situation itself must NOT be declared as a pointer, because its label is defined inside the retained context and declaring it twice duplicates the label; the first revision did that and the engine refused the file with a multiply-defined-label warning. The situation is context rather than a counted unit and does not overlap the counted interval.\par\medskip
\noindent\textit{Editorial note.} The retained proof names four results of this chapter and two results of the Varieties chapter. All are declared as non-load-bearing navigation pointers with their stable Stacks tags looked up in the shared tag index; the chapter ones are counted by cards of this lane and none of their intervals is reproduced here.\par\medskip
\noindent\textit{Editorial note.} The retained proof span starts at line 5232, the line of the opening marker read from the source, and contains exactly one proof environment.\par\medskip
\noindent\textit{Editorial note.} No citation occurs in the retained spans, so no bibliography entry is supplied.\par\medskip
\noindent\textit{Allowed standard prerequisites (editorial).}\begin{itemize}
\item The regular-model situation, the five defining conditions of a numerical type and the genus formula. The situation fixes a discrete valuation ring with residue field, a smooth projective curve over its fraction field, a regular proper model and the components of the special fibre with their multiplicities; the numerical type attached to the model records the number of components, their multiplicities, their intersection numbers, the degrees of their residue fields over the residue field of the ring and their genera, and the five defining conditions of a numerical type are positivity, symmetry with nonnegative off-diagonal entries, connectedness of the intersection graph, vanishing of the weighted row sums and divisibility of the intersection numbers by the weights. The genus of the curve is given by the chapter's formula in terms of that data, and the intersection numbers are computed by the chapter's degree pairing together with the divisibility of degrees by the residue field degree. The author's own definitions and lemmas of the chapter and the quoted result of the Varieties chapter are the allowed inputs.. The counted lemma is the bridge between the geometry of a regular model and the combinatorics of numerical types, and every later statement of the section is phrased in terms of the type it produces.
\end{itemize}\medskip
\ifdefined\hypertarget\hypertarget{sourcepointer-1}{}\fi
\noindent\textit{Source reference.} Original definition 0C6Z: focused-seeds/AGP/sources/stacks-snapshot-20261011/tags.txt, source lines 15314--15314. Navigation pointer to the chapter's definition of a numerical type with its five conditions, which the retained proof verifies one by one; the cited definition is an allowed input and its text is not reproduced here. The number is the stable Stacks tag of the target, looked up in the shared tag index.\par\medskip
\ifdefined\hypertarget\hypertarget{sourcepointer-2}{}\fi
\noindent\textit{Source reference.} Original definition 0C72: focused-seeds/AGP/sources/stacks-snapshot-20261011/tags.txt, source lines 15317--15317. Navigation pointer to the chapter's definition of the genus of a numerical type; the cited definition is an allowed input and its text is not reproduced here. The number is the stable Stacks tag of the target, looked up in the shared tag index.\par\medskip
\ifdefined\hypertarget\hypertarget{sourcepointer-3}{}\fi
\noindent\textit{Source reference.} Original equation 0C65: focused-seeds/AGP/sources/stacks-snapshot-20261011/tags.txt, source lines 15285--15285. Navigation pointer to the chapter's displayed formula for the intersection numbers, used by the retained proof when it computes the third entry of the tuple; the cited labelled equation is an allowed input and its text is not reproduced here. The number is the stable Stacks tag of the target, looked up in the shared tag index.\par\medskip
\ifdefined\hypertarget\hypertarget{sourcepointer-4}{}\fi
\noindent\textit{Source reference.} Original lemma 0C64: focused-seeds/AGP/sources/stacks-snapshot-20261011/tags.txt, source lines 15284--15284. Navigation pointer invoked by the retained proof for the symmetry of the degree pairing, which gives the symmetry of the intersection matrix; that lemma is counted by a separate card of this lane and its interval is not reproduced here. The number is the stable Stacks tag of the target, looked up in the shared tag index.\par\medskip
\ifdefined\hypertarget\hypertarget{sourcepointer-5}{}\fi
\noindent\textit{Source reference.} Original lemma 0C66: focused-seeds/AGP/sources/stacks-snapshot-20261011/tags.txt, source lines 15286--15286. Navigation pointer invoked by the retained proof for the numerical properties of the intersection form; the cited statement is an allowed input of the same chapter and its proof is not reproduced here. The number is the stable Stacks tag of the target, looked up in the shared tag index.\par\medskip
\ifdefined\hypertarget\hypertarget{sourcepointer-6}{}\fi
\noindent\textit{Source reference.} Original lemma 0CA3: focused-seeds/AGP/sources/stacks-snapshot-20261011/tags.txt, source lines 15423--15423. Navigation pointer invoked by the retained proof for the formula expressing the genus of the curve through the data of the numerical type; that lemma is counted by a separate card of this lane and its interval is not reproduced here. The number is the stable Stacks tag of the target, looked up in the shared tag index.\par\medskip
\ifdefined\hypertarget\hypertarget{sourcepointer-7}{}\fi
\noindent\textit{Source reference.} Original lemma 0C6P: focused-seeds/AGP/sources/stacks-snapshot-20261011/tags.txt, source lines 15304--15304. Allowed input of the Varieties chapter quoted by the retained proof for the divisibility of the degree of an invertible sheaf on a component by its residue field degree, which gives the divisibility condition of the type; the pointer is registered so the reference resolves and the cited statement is not reproduced inside the span. The number is the stable Stacks tag of the target, looked up in the shared tag index.\par\medskip
\ifdefined\hypertarget\hypertarget{sourcepointer-8}{}\fi
\noindent\textit{Source reference.} Original lemma 0C5Z: focused-seeds/AGP/sources/stacks-snapshot-20261011/tags.txt, source lines 15279--15279. Navigation pointer carried by the retained situation and invoked by the retained proof for the vanishing of the weighted row sums, which is the statement that the special fibre is a principal divisor; that lemma is counted by a separate card of this lane and its interval is not reproduced here. The number is the stable Stacks tag of the target, looked up in the shared tag index.\par\medskip
\begin{situation}
\label{situation-regular-model}
Let $R$ be a discrete valuation ring with fraction field $K$,
residue field $k$, and uniformizer $\pi$.
Let $C$ be a smooth projective curve over $K$ with $H^0(C, \mathcal{O}_C) = K$.
Let $X$ be a regular proper model of $C$.
Let $C_1, \ldots, C_n$ be the irreducible components of the special
fibre $X_k$. Write $X_k = \sum m_i C_i$ as in
Lemma \ref{lemma-divisor-special-fiber}.
\end{situation}

\begin{lemma}
\label{lemma-numerical-type-of-model}
In Situation \ref{situation-regular-model} with
$\kappa_i = H^0(C_i, \mathcal{O}_{C_i})$ and $g_i$ the genus of $C_i$
the data
$$
n, m_i, (C_i \cdot C_j), [\kappa_i : k], g_i
$$
is a numerical type of genus equal to the genus of $C$.
\end{lemma}

\begin{proof}
(In the proof of Lemma \ref{lemma-genus-formula}
we have seen that the quantities
used in the statement of the lemma are well defined.)
We have to verify the conditions (1) -- (5) of
Definition \ref{definition-type}.

\medskip\noindent
Condition (1) is immediate.

\medskip\noindent
Condition (2). Symmetry of the matrix $(C_i \cdot C_j)$ follows from
Equation (\ref{equation-form}) and
Lemma \ref{lemma-intersection-pairing}.
Nonnegativity of $(C_i \cdot C_j)$ for $i \not = j$
is part (1) of Lemma \ref{lemma-properties-form}.

\medskip\noindent
Condition (3) is part (3) of Lemma \ref{lemma-properties-form}.

\medskip\noindent
Condition (4) is part (2) of Lemma \ref{lemma-properties-form}.

\medskip\noindent
Condition (5) follows from the fact that $(C_i \cdot C_j)$ is
the degree of an invertible module on $C_i$ which is divisible
by $[\kappa_i : k]$, see Varieties, Lemma \ref{varieties-lemma-divisible}.

\medskip\noindent
The genus formula proved in Lemma \ref{lemma-genus-formula}
tells us that the numerical type has the genus as stated, see
Definition \ref{definition-genus}.
\end{proof}

\end{document}
